%% BioCompiler (VSTTE submission) — Sections 1–3
%% "Well-Typed Genes Don't Go Wrong: A Type System for Machine-Verified Gene Design"
%%
%% This file contains Sections 1–3 as a self-contained LaTeX fragment.
%% Compile within the main document that defines all custom commands.
%%
%% Required packages (must be in preamble of main document):
%%   mathpartir, booktabs, enumitem, xspace, hyperref, tikz,
%%   amsmath, amssymb, amsfonts, xcolor
%%
%% Required custom commands (defined in main document preamble):
%%   \PASS, \FAIL, \UNCERTAIN, \tvl, \tand, \seq, \pred,
%%   \verdict, \judg, \GOLD, \SILVER, \BRONZE

% =================================================================
\section{Introduction}\label{sec:intro}

Synthetic biology promises custom-designed proteins for medicine,
industry, and research, yet gene design today relies on heuristics
without formal guarantees.
A designed gene may harbor cryptic splice sites, out-of-frame exons,
protein misfolding risks, or immunogenic epitopes---flaws discovered
only after costly wet-lab experiments, failed clinical trials, or, in
the worst case, adverse patient outcomes.
The obstacle appears fundamental: biology is non-deterministic, so how
can one \emph{prove} a property of a probabilistic system?

\paragraph{Why constraint lists fail.}
Existing gene design tools---DNA~Chisel~\cite{piret2020dnachisel},
GeneDesign~\cite{richardson2006genedesign},
DNAworks~\cite{vandenbergh2015dnaworks}---frame the problem as
\emph{constraint satisfaction}: provide a list of constraints, run an
optimizer, check whether each constraint passes.
This approach has three fundamental weaknesses:

\begin{enumerate}[leftmargin=*,itemsep=1pt,topsep=2pt]
\item \textbf{Constraints don't compose.}
Satisfying constraint~$C_1$ may violate~$C_2$ in a previously
processed region, or create violations at fragment junctions.
DNA~Chisel applies solving passes sequentially; a later pass can undo
the work of an earlier one, with no global soundness guarantee.
The problem is not merely theoretical: we present a concrete
counterexample where two individually passing gene fragments produce a
Sau3AI restriction site (\texttt{GATC}) at their junction
(\S\ref{sec:motivating}).

\item \textbf{No certificates.}
A constraint-based tool reports \texttt{PASS} or \texttt{FAIL} for
each constraint, but provides no independently verifiable proof.
The entire optimizer is in the trusted computing base (TCB): a bug in
any solving pass can produce a false \texttt{PASS}, and there is no
mechanism for an independent auditor to detect this without re-running
the full toolchain.

\item \textbf{Binary logic cannot handle uncertainty.}
Binary constraint satisfaction offers two outcomes:
\texttt{PASS} or \texttt{FAIL}.
But gene design involves inherent uncertainty: a splice site scores
2.5~bits, above the noise floor but below the definitive threshold.
A binary system must either round up to \FAIL{} (over-conservative,
wasting potentially viable designs) or round down to \PASS{}
(unsound, potentially approving a flawed design).
The problem is especially acute for valine codons, which
\emph{invariably} contain the GT dinucleotide---the universal
5\textquotesingle{} splice donor signal---making binary classification
at valine positions fundamentally inadequate.
\end{enumerate}

\paragraph{Our approach.}
We recast gene design as a \emph{type-checking} problem.
Each biological requirement becomes a type predicate; the nucleotide
sequence is the term; and a five-valued logic yields
verdicts $\{\PASS, \LIKELYPASS, \UNCERTAIN, \LIKELYFAIL, \FAIL\}$,
with a three-valued subset (Kleene~K3: $\{\PASS, \UNCERTAIN, \FAIL\}$)
formalized in Lean~4.
The key insight:

\begin{quote}
\emph{Types compose; constraints don't.}
\end{quote}

\noindent Under a type system, combining predicates via Kleene
conjunction ($\sqcap$) is a \emph{theorem}, not an engineering
obligation:
\[
  \judg{\Gamma}{\seq}{\pred_1 \sqcap \pred_2}
  \;=\;
  \judg{\Gamma}{\seq}{\pred_1} \;\sqcap\;
  \judg{\Gamma}{\seq}{\pred_2}
\]
If the composite verdict is \PASS, then \emph{each individual property
holds}. No manual reasoning about interactions is needed. This is the
same algebraic guarantee that type systems provide in programming
languages: knowing that $f : A \to B$ and $x : A$ lets you conclude
$f(x) : B$ without inspecting the implementation.

Three-valued logic avoids the false dichotomy of binary constraint
satisfaction: \PASS{} means the property is unconditionally guaranteed;
\FAIL{} means it is guaranteed \emph{not} to hold; \UNCERTAIN{} means
the system declines to guarantee, which is the honest answer when
evidence is insufficient.
The soundness theorem guarantees that \PASS{} is never wrong.

\paragraph{The CompCert analogy.}
Like CompCert~\cite{leroy2009compcert}, BioCompiler separates the
\emph{unverified} optimizer from the \emph{verified} type checker.
A bug in the optimizer can cause a \FAIL{} verdict (a false negative)
but never a false \PASS.
This separation is not merely architectural: it is \emph{proved} as a
theorem (SLOT-Independence, Corollary~\ref{cor:slot}), which
guarantees that the type checker's verdicts are independent of the
optimizer's internal state.
A standalone verifier (${\sim}460$~LOC, zero biocompiler dependencies)
further reduces the TCB: certificate consumers need only trust the
verifier plus the Lean~4 kernel.

\paragraph{SLOT: Sealed Local Oracle Theory.}
A key challenge in applying formal methods to gene design is that many
biological properties---protein structure, stability, immunogenicity---
require external analysis tools (ESMFold, FoldX, CamSol, NetMHCpan)
that cannot be formalized within Lean~4.
We introduce \emph{Sealed Local Oracle Theory} (SLOT), a framework
that treats external tool outputs as \emph{sealed oracles}: their
results are used to inform optimization decisions but are
\emph{explicitly excluded} from the type-checking judgment.
SLOT operates in three modes:
\emph{CONSERVATIVE} (no external tool may be consulted; only
computationally decidable predicates are evaluated),
\emph{VERIFIED} (external tools may be consulted, but their outputs
are not trusted for \PASS{} verdicts), and
\emph{PERMISSIVE} (external tool outputs may influence verdicts, at
the cost of extending the TCB).
We prove a formal \emph{refinement theorem}:
\emph{VERIFIED} mode refines \emph{CONSERVATIVE} mode, meaning that
every \PASS{} verdict in VERIFIED mode is also valid in CONSERVATIVE
mode.
This gives users a principled knob: turn up the trust level to gain
coverage, with a formal guarantee that the result is at least as strong
as the conservative baseline.

\paragraph{Contributions.}
\begin{enumerate}[leftmargin=*,itemsep=1pt,topsep=2pt]
\item A five-valued type system for gene design (three-valued Kleene~K3
      subset formalized in Lean~4)---\emph{the first machine-verified type
      system for codon-level gene design constraints}---proved in
      Lean~4 across 17~modules (12 proof $+$ 5 library) and
      8,495~lines of proof code, with zero \texttt{sorry} in the core (15 in SLOTVerification.lean) and zero standalone
      axioms (\S\ref{sec:typesys}).  We acknowledge BBN's
      Proto-BioCompiler~\cite{beal2011protobiocompiler} as early prior
      art that applied compiler concepts to biological network design;
      BioCompiler differs by targeting codon-level gene design with
      machine-checked soundness proofs, whereas Proto-BioCompiler's
      correctness arguments were informal.

\item Compositional soundness: combining predicates via Kleene
      conjunction preserves guarantees, proved as a theorem.
      We demonstrate that constraint lists fail composition via a
      concrete junction counterexample (Sau3AI at a GATC boundary)
      (\S\ref{sec:motivating}, \S\ref{sec:compositionality}).

\item Forty-three type predicates spanning five biological domains
      (DNA-level, structure, stability, solubility, immunogenicity),
      with formal three-valued semantics for each
      (\S\ref{sec:predicates}).

\item The SLOT (Sealed Local Oracle Theory) framework for
      integrating external analysis tools without extending the TCB,
      with a proved refinement theorem establishing that VERIFIED mode
      refines CONSERVATIVE mode (\S\ref{sec:slot}).

\item Graduated certificates (\GOLD/\SILVER/\BRONZE) communicating
      design confidence, with a type-directed protein mutagenesis
      mechanism using BLOSUM62-guided amino acid substitution for
      resolving structurally unavoidable violations, and a standalone
      verifier (${\sim}460$~LOC, zero biocompiler dependencies)
      enabling independent certificate re-checking.

\item Retrospective validation against four published clinical and
      experimental datasets (SCID gene therapy, $\beta$-thalassemia,
      $\alpha$-synuclein aggregation, EPO immunogenicity), plus
      property-based testing with 500~Hypothesis tests bridging the
      Lean~4 model to the Python implementation, across a total of
      13,359~tests comprising ~169,000~test lines.
\end{enumerate}

% =================================================================
\section{Motivating Example}\label{sec:motivating}

We illustrate BioCompiler's advantages through a concrete gene design
scenario: synthesizing human insulin for recombinant expression in
\emph{E.~coli}.
Insulin is an ideal motivating example because it is short enough to
reason about manually (preproinsulin, 86~aa), yet involves all the
compositional challenges that arise in production gene design.

\subsection{The Problem: Insulin Gene Design}

Human preproinsulin requires a coding sequence that simultaneously
satisfies multiple biological requirements spanning all five predicate
domains:
\begin{itemize}[leftmargin=*,itemsep=0pt,topsep=1pt]
\item \emph{DNA-level}: no cryptic splice sites, no restriction
      enzyme recognition sequences, codon adaptation index $\ge 0.8$
      for the expression host, GC content in $[30\%, 70\%]$, reading
      frame maintained, no instability motifs, no CpG islands.
\item \emph{Structure}: the designed sequence must not introduce
      misfolding risk or unexpected transmembrane domains.
\item \emph{Stability}: the folded protein must have stable free
      energy, intact disulfide bonds, and a well-packed hydrophobic
      core.
\item \emph{Solubility}: the protein must express in soluble form,
      without aggregation-prone regions or excessive hydrophobic
      stretches.
\item \emph{Immunogenicity}: the designed protein should not contain
      strong T-cell epitopes or dominant B-cell epitopes for the
      target population.
\end{itemize}

\noindent In a conventional constraint-based tool, each requirement is
an independent entry in a checklist.
The tool iterates through the list, attempting to satisfy each
constraint in turn.
This approach breaks down in two ways that we now illustrate.

\subsection{Fragment Composition Creates Junction Violations}

\begin{figure}[t]
\centering
\begin{tikzpicture}[
  frag/.style={draw,rounded corners,minimum height=0.8cm,
               minimum width=3.2cm,font=\small,align=center},
  arr/.style={-{Stealth[length=5pt]},thick},
  note/.style={font=\scriptsize,text width=3.0cm,align=left}
]
% Fragment A
\node[frag,fill=green!8] (fa) {Fragment A: \dots\texttt{GAT}};
\node[note,right=0.3cm of fa] (na) {Last codon: GAT (Asp)\\
  No \texttt{GATC} internally\\$\Rightarrow$ \PASS\ $\checkmark$};

% Fragment B
\node[frag,fill=green!8,below=0.8cm of fa] (fb) {Fragment B: \texttt{CTT}\dots};
\node[note,right=0.3cm of fb] (nb) {First codon: CTT (Leu)\\
  No \texttt{GATC} internally\\$\Rightarrow$ \PASS\ $\checkmark$};

% Concatenation
\node[frag,fill=red!12,below=0.8cm of fb] (ab) {A$+$B: \dots\texttt{GATC}TT\dots};
\node[note,right=0.3cm of ab] (nab) {Junction: \texttt{GATC}\\
  Sau3AI site.\\$\Rightarrow$ \FAIL\ $\times$};

\draw[arr] (fa) -- node[left,font=\scriptsize] {concatenate} (fb);
\draw[arr] (fb) -- node[left,font=\scriptsize] {concatenate} (ab);

% Highlight junction
\draw[red,thick,decorate,decoration={brace,amplitude=5pt,mirror}]
  ([xshift=1.2cm]ab.north west) -- ([xshift=2.2cm]ab.north west)
  node[midway,below=6pt,font=\scriptsize,red] {Sau3AI};
\end{tikzpicture}
\caption{Compositional failure: two individually passing fragments
         create a Sau3AI restriction site at their junction.
         Fragment~A ends with codon GAT (Aspartic acid); Fragment~B
         begins with CTT (Leucine). Neither contains \texttt{GATC}
         internally, but their concatenation does.}
\label{fig:junction}
\end{figure}

Consider designing insulin by assembling two gene fragments, each
checked independently against the \emph{NoRestrictionSite} constraint.
Figure~\ref{fig:junction} illustrates the scenario.
Fragment~A ends with the codon GAT (encoding Aspartic acid, D);
Fragment~B begins with CTT (encoding Leucine, L).
Neither fragment contains the Sau3AI recognition sequence
\texttt{GATC} internally, so each passes the constraint check
independently.
However, the concatenation $A \cdot B$ creates the sequence
\dots\texttt{GATC}TT\dots, where the \texttt{GATC} straddles the
junction---the T from GAT and the C from CTT together form the
restriction site.

\begin{table}[t]
\centering
\caption{Compositional failure: two individually passing fragments
         create a Sau3AI site at their junction.}
\label{tab:junction}
\begin{tabular}{llll}
\toprule
\textbf{Fragment} & \textbf{Ends} & \textbf{Last/First Codon} &
  \textbf{Sau3AI-free?} \\
\midrule
A & \dots\texttt{GAT} & GAT (Asp, D) & $\checkmark$ \PASS \\
B & \texttt{CTT}\dots & CTT (Leu, L) & $\checkmark$ \PASS \\
A+B & \dots\texttt{GATC}TT\dots & Junction: \texttt{GATC} &
  $\times$ \FAIL \\
\bottomrule
\end{tabular}
\end{table}

A checklist approach processes each fragment independently:
\begin{enumerate}[leftmargin=*,itemsep=0pt,topsep=1pt]
\item Check Fragment~A against [\emph{no restriction sites}]
      $\to$ \PASS~$\checkmark$
\item Check Fragment~B against [\emph{no restriction sites}]
      $\to$ \PASS~$\checkmark$
\item Concatenate A $+$ B
\item No re-check step. $\to$ \texttt{GATC} exists but was never
      detected
\end{enumerate}

The fundamental issue is that checklist items are checked
independently.
There is no compositional logic---``\PASS{} $\wedge$ \PASS{}
$\to$ \PASS'' is \emph{assumed} but \emph{not guaranteed}.
Junction regions are invisible to per-fragment checking.
This is not a theoretical curiosity: gene synthesis companies
routinely assemble oligonucleotide fragments, and junction
restriction sites are a known class of design failure
\cite{piret2020dnachisel}.

BioCompiler's type system enforces re-verification on the composed
artifact: the cross-codon restriction scanner
\texttt{find\_cross\_codon\_restriction()} detects the cross-boundary
\texttt{GATC} site.
The fix is trivial: change Fragment~A's last codon from GAT to GAC
(both encode Aspartic acid).
The junction becomes \texttt{GACCTT}---no Sau3AI site.
The type system's compositional soundness theorem
(Corollary~\ref{cor:comp}) guarantees that such junction violations
cannot escape detection: if the composite verdict is \PASS, then
\emph{every} predicate holds on the \emph{entire} composed sequence,
including junctions.

\subsection{UNCERTAIN: Valine Codons Inherently Contain GT}

Not every borderline splice site warrants a definitive \FAIL.
Consider positions where valine is required in the insulin protein.
All four valine codons (GTT, GTC, GTA, GTG) contain the GT
dinucleotide---the universal 5\textquotesingle{} splice donor signal.
This is not a deficiency of the codon table; it is a fundamental
property of the genetic code.
A binary constraint system faces an impossible choice:
\begin{itemize}[leftmargin=*,itemsep=0pt,topsep=1pt]
\item Round up to \FAIL{} (over-conservative: every valine position
      in the gene fails the cryptic splice check, even though most
      valine positions are splicing-neutral in context), or
\item Round down to \PASS{} (unsound: a genuine cryptic donor at a
      valine position is silently approved).
\end{itemize}

Neither option is correct.
The over-conservative choice eliminates viable designs unnecessarily;
the unsound choice approves designs that may cause aberrant splicing
in the patient.
Both outcomes have clinical precedent: over-conservative filtering
delays therapeutic development, while unsound approval has contributed
to adverse events in gene therapy trials~\cite{hacein2018gene}.

BioCompiler's dual-threshold scheme provides a third option:
\begin{itemize}[leftmargin=*,itemsep=0pt,topsep=1pt]
\item \PASS{} if the MaxEntScan score $< \theta_u = 1.5$~bits
      (clearly safe---the GT dinucleotide exists but the surrounding
      context is sufficiently non-spliceogenic),
\item \UNCERTAIN{} if $\theta_u \le \mathit{score} < \theta_c$
      (borderline---the site is neither clearly safe nor clearly
      dangerous),
\item \FAIL{} if $\mathit{score} \ge \theta_c = 3.0$~bits
      (clearly unsafe---the surrounding context makes the GT
      dinucleotide a strong cryptic donor).
\end{itemize}

This is not a cop-out---it is a \emph{precise characterization} of
the state of knowledge.
Valine codons inherently contain GT, so the type system correctly
reports \UNCERTAIN{} at positions where the GT falls in a borderline
scoring context, rather than a misleading \PASS{} or an
over-conservative \FAIL.
The soundness theorem guarantees that \PASS{} is never wrong;
\UNCERTAIN{} means the system declines to guarantee.

The dual-threshold approach mirrors established practice in medical
diagnostics and risk assessment, where borderline results are reported
as ``indeterminate'' rather than forced into a binary classification.
In the context of gene design, this third verdict drives a principled
escalation: \UNCERTAIN{} positions are candidates for type-directed
mutagenesis (\S\ref{sec:mut}), where conservative amino acid
substitutions (e.g., valine $\to$ isoleucine, BLOSUM62~$= +3$) may
eliminate the problematic motif entirely.

\subsection{Certificate Trust and Graduated Guarantees}

After optimization and type checking, BioCompiler emits a graduated
certificate:
\begin{description}[leftmargin=!,labelwidth=1.4cm,itemsep=1pt,
                     topsep=1pt]
\item[\GOLD] All predicates \PASS{} through codon optimization
      alone---no protein modification required.
      The designed nucleotide sequence satisfies every predicate
      without altering the target amino acid sequence.
\item[\SILVER] \PASS{} achieved after type-directed mutagenesis---
      conservative amino acid substitutions were necessary at positions
      where all synonymous codons violate a predicate (e.g.,
      GT-mandatory valine positions).
      Protein identity is preserved above 97\%.
\item[\BRONZE] One or more predicates remain \UNCERTAIN{}
      after exhausting all conservative substitutions.
      The certificate documents which predicates are uncertain and why.
\end{description}

The certificate is not merely a label---it is a \emph{machine-readable
record} containing the sequence, each predicate's verdict with
derivation traces, and provenance metadata (tool version, proof commit
hash, SLOT mode, timestamp).
An independent verifier (${\sim}460$~LOC, importing only the Python
standard library) re-checks every predicate from scratch, using only
the sequence and parameters embedded in the certificate.
The optimizer is \emph{outside} the TCB: even a buggy optimizer
cannot produce a false \PASS{} certificate, because the verifier
re-evaluates all predicates independently.

The graduated guarantee system provides a principled communication
channel between the gene design tool and the downstream
biologist.
A \GOLD{} certificate says: ``this design is guaranteed correct at the
nucleotide level.''
A \SILVER{} certificate says: ``this design is guaranteed correct
after conservative protein modifications; here are the substitutions
that were made.''
A \BRONZE{} certificate says: ``this design has unresolved concerns;
human judgment is required.''
This graduation is essential for practical adoption: biologists are
accustomed to making trade-offs (e.g., accepting a slightly less
optimal codon adaptation index in exchange for eliminating a cryptic
splice site), and the certificate system makes these trade-offs
\emph{explicit} and \emph{auditable}.

% =================================================================
\section{The Type System}\label{sec:typesys}

We formalize BioCompiler's type system, which assigns three-valued
verdicts to nucleotide sequences relative to biological predicates
spanning five domains: DNA-level, structure, stability, solubility,
and immunogenicity.
The formalization consists of the verdict algebra (Kleene~K3),
forty-three type predicates (thirty-six with three-valued semantics in Lean~4; seven additional extended diagnostic predicates in the Python implementation), typing
rules, and soundness/compositionality theorems---all machine-verified
in Lean~4.

\subsection{Three-Valued Logic (Kleene K3)}\label{sec:3vl}

The set of verdicts is $\tvl = \{\PASS, \FAIL, \UNCERTAIN\}$.
We equip $\tvl$ with two operations: Kleene conjunction ($\sqcap$) and
Kleene disjunction ($\sqcup$), following Kleene's strong three-valued
logic~K3~\cite{kleene1952introduction}.

\begin{figure}[t]
\centering
\begin{tabular}{c|ccc}
$\sqcap$ & \PASS & \UNCERTAIN & \FAIL \\ \hline
\PASS     & \PASS     & \UNCERTAIN     & \FAIL \\
\UNCERTAIN & \UNCERTAIN & \UNCERTAIN     & \FAIL \\
\FAIL     & \FAIL     & \FAIL          & \FAIL
\end{tabular}
\hfill
\begin{tabular}{c|ccc}
$\sqcup$ & \PASS & \UNCERTAIN & \FAIL \\ \hline
\PASS     & \PASS     & \PASS          & \PASS \\
\UNCERTAIN & \PASS     & \UNCERTAIN     & \UNCERTAIN \\
\FAIL     & \PASS     & \UNCERTAIN     & \FAIL
\end{tabular}
\caption{Kleene K3 truth tables for conjunction ($\sqcap$) and
         disjunction ($\sqcup$). \FAIL{} is absorptive for $\sqcap$;
         \PASS{} is absorptive for $\sqcup$. The key property for
         soundness is: $v_1 \sqcap v_2 = \PASS \Rightarrow v_1 =
         v_2 = \PASS$.}
\label{fig:truthtables}
\end{figure}

The choice of Kleene K3 over other three-valued logics ({\L}ukasiewicz
\L$_3$, Bochvar~B$_3$, Priest~LP) is deliberate and motivated by the
semantics of gene design:

\begin{itemize}[leftmargin=*,itemsep=0pt,topsep=1pt]
\item \textbf{Absorptivity of \FAIL}:
$v_1 \sqcap \FAIL = \FAIL$ for all~$v_1$.
If any predicate is definitively violated, the composite design is
violated, regardless of the state of other predicates.
This matches the safety-critical nature of gene design: a single
guaranteed violation suffices to reject the design.

\item \textbf{Soundness-preserving conjunction}:
$v_1 \sqcap v_2 = \PASS \Rightarrow v_1 = v_2 = \PASS$.
This is the algebraic backbone of compositional soundness: a composite
\PASS{} verdict decomposes into individual \PASS{} verdicts, each of
which is backed by the per-predicate soundness theorem.

\item \textbf{Uncertainty is contagious but not fatal}:
$v_1 \sqcap \UNCERTAIN = \UNCERTAIN$ when $v_1 \neq \FAIL$.
A single uncertain predicate downgrades the composite from \PASS{} to
\UNCERTAIN{} but does not produce a \FAIL.
This matches the biological intuition: uncertainty weakens confidence
but does not establish a violation.
\end{itemize}

\noindent Kleene K3 forms a \emph{commutative semilattice} under
$\sqcap$ (with \PASS{} as the top element and \FAIL{} as the bottom),
ensuring associativity, commutativity, and idempotency---all of which
are proved in Lean~4 (module \texttt{ThreeValued.lean}).

\subsection{Type Predicates}\label{sec:predicates}

Let $\Sigma = \{\texttt{A}, \texttt{C}, \texttt{G}, \texttt{T}\}$
denote the DNA alphabet.
A \emph{DNA sequence} $\seq \in \Sigma^*$ is a finite string over
$\Sigma$.
A \emph{typing context} $\Gamma$ maps sequence positions to biological
annotations (exon boundaries, reading frame, organism, cell type,
target population HLA allele frequencies).

BioCompiler defines \textbf{28~type predicates} organized in five
domains, each mapping a (predicate, sequence, context) triple to a
verdict in~$\tvl$.

\subsubsection{DNA-Level Predicates (12)}

The DNA-level predicates operate directly on the nucleotide sequence
and are computationally decidable without external tools.
They form the \emph{core} of the type system: their verdicts depend
only on the sequence and parameters, with no SLOT involvement.

\begin{table}[t]
\centering
\caption{DNA-level type predicates. All are computationally decidable
         from the sequence alone; none require SLOT oracles.
         Predicates marked with $\dagger$ use dual-threshold
         semantics.}
\label{tab:dna-predicates}
\begin{tabular}{llp{4.6cm}}
\toprule
\textbf{Predicate} & \textbf{Parameters} & \textbf{Three-valued semantics} \\
\midrule
$\mathit{NoStopCodons}$ & --- &
  \PASS{} iff no in-frame stop codon; \FAIL{} otherwise \\
$\mathit{SpliceCorrect}(C)$ & Exon--intron structure &
  \PASS{} iff NDFST produces unique isoform matching $C$ \\
$\mathit{NoCrypticSplice}(\theta_c,\theta_u)^\dagger$ & Dual thresholds &
  \PASS{} if all sites $< \theta_u$; \UNCERTAIN{} if $\in
  [\theta_u, \theta_c)$; \FAIL{} if $\ge \theta_c$ \\
$\mathit{NoCrypticPromoter}(\theta_c,\theta_u)^\dagger$ & Dual thresholds &
  \PASS{} if no promoter-like motif above $\theta_u$; \UNCERTAIN{}
  if borderline; \FAIL{} if above $\theta_c$ \\
$\mathit{CodonAdapted}(O,\theta)$ & Organism, threshold &
  \PASS{} iff $\mathit{CAI}(\seq, O) \ge \theta$ \\
$\mathit{GCInRange}(lo,hi)$ & Bounds &
  \PASS{} iff $lo \le \mathit{gc}(\seq) \le hi$ \\
$\mathit{NoRestrictionSite}(S)$ & Enzyme set &
  \PASS{} iff no site from $S$ occurs in $\seq$ (incl.\ junctions) \\
$\mathit{InFrame}(rf,b)$ & Frame, base &
  \PASS{} iff frame consistent, no premature stop \\
$\mathit{NoInstabilityMotif}$ & --- &
  \PASS{} iff no ATTTA or U-rich degradation motif \\
$\mathit{NoCpGIsland}$ & --- &
  \PASS{} iff no 200\,bp window with Obs/Exp $> 0.6$ \\
$\mathit{NoUnexpectedTMDomain}(\theta)^\dagger$ & Threshold &
  \PASS{} if no TM segment $> \theta$~aa; \UNCERTAIN{} if
  borderline \\
$\mathit{CodonOptimality}(O,\theta)$ & Organism, threshold &
  \PASS{} iff codon optimality score $\ge \theta$ \\
\bottomrule
\end{tabular}
\end{table}

\noindent The DNA-level predicates capture the classical concerns of
gene design: reading frame integrity, splice fidelity, codon usage
adaptation, GC content, restriction site avoidance, mRNA stability,
CpG methylation risk, and transmembrane domain prediction.
With the exception of $\mathit{NoUnexpectedTMDomain}$ (which uses a
heuristic hydropathy sliding window), all DNA-level predicates are
decidable by direct computation on the sequence.
The dual-threshold predicates ($\mathit{NoCrypticSplice}$,
$\mathit{NoCrypticPromoter}$, $\mathit{NoUnexpectedTMDomain}$) are
the primary sources of \UNCERTAIN{} verdicts: they classify borderline
signals that cannot be safely committed to either \PASS{} or \FAIL{}.

\subsubsection{Structure Predicates (4)}

Structure predicates assess whether the designed sequence will fold
into the intended three-dimensional structure.
These predicates require external structure prediction tools and are
\emph{SLOT predicates}: their verdicts in CONSERVATIVE mode are always
\UNCERTAIN{} (the type system cannot verify structural properties
without oracle input), while in VERIFIED mode they may yield \PASS{}
or \FAIL{} based on oracle output---but the soundness guarantee is
maintained by the SLOT-Independence theorem.

\begin{table}[t]
\centering
\caption{Structure type predicates. All are SLOT predicates; verdicts
         depend on external structure prediction tools (ESMFold
         API with heuristic fallback). In CONSERVATIVE mode, all
         yield \UNCERTAIN.}
\label{tab:struct-predicates}
\begin{tabular}{llp{4.6cm}}
\toprule
\textbf{Predicate} & \textbf{Parameters} & \textbf{Three-valued semantics} \\
\midrule
$\mathit{StructureConfidence}(\theta)^\dagger$ & Threshold &
  \PASS{} if pLDDT $\ge \theta$; \UNCERTAIN{} if borderline;
  \FAIL{} if pLDDT $< \theta_{\mathit{lo}}$ \\
$\mathit{NoMisfoldingRisk}(\theta)$ & Threshold &
  \PASS{} if predicted TM-score vs.\ wildtype $\ge \theta$;
  \FAIL{} otherwise \\
$\mathit{CorrectFoldTopology}$ & --- &
  \PASS{} if predicted fold matches expected topology class \\
$\mathit{NoUnexpectedInteraction}$ & --- &
  \PASS{} if no predicted intermolecular interface \\
\bottomrule
\end{tabular}
\end{table}

\subsubsection{Stability Predicates (4)}

Stability predicates evaluate the thermodynamic stability of the folded
protein, using FoldX (CLI with heuristic fallback) as the analysis
engine.

\begin{table}[t]
\centering
\caption{Stability type predicates. All are SLOT predicates; verdicts
         depend on FoldX stability analysis.}
\label{tab:stab-predicates}
\begin{tabular}{llp{4.6cm}}
\toprule
\textbf{Predicate} & \textbf{Parameters} & \textbf{Three-valued semantics} \\
\midrule
$\mathit{StableFolding}(\theta)^\dagger$ & $\Delta G$ threshold &
  \PASS{} if $\Delta G < \theta$; \UNCERTAIN{} if borderline;
  \FAIL{} if $\Delta G \ge \theta_{\mathit{hi}}$ \\
$\mathit{NoDestabilizingMutation}$ & --- &
  \PASS{} if $\Delta\Delta G < 1.0$~kcal/mol for all mutations \\
$\mathit{DisulfideBondIntegrity}$ & --- &
  \PASS{} if all annotated disulfide bonds are intact in prediction \\
$\mathit{HydrophobicCoreQuality}(\theta)$ & Threshold &
  \PASS{} if core packing score $\ge \theta$ \\
\bottomrule
\end{tabular}
\end{table}

\subsubsection{Solubility Predicates (4)}

Solubility predicates assess whether the protein will express in
soluble form, using CamSol (Wimley-White and Urry hydrophobicity
scales) as the analysis engine.

\begin{table}[t]
\centering
\caption{Solubility type predicates. All are SLOT predicates; verdicts
         depend on CamSol solubility analysis.}
\label{tab:sol-predicates}
\begin{tabular}{llp{4.6cm}}
\toprule
\textbf{Predicate} & \textbf{Parameters} & \textbf{Three-valued semantics} \\
\midrule
$\mathit{SolubleExpression}(\theta)^\dagger$ & Threshold &
  \PASS{} if CamSol score $\ge \theta$; \UNCERTAIN{} if
  borderline; \FAIL{} if $< \theta_{\mathit{lo}}$ \\
$\mathit{NoAggregationProneRegion}(w)$ & Window size &
  \PASS{} if no window of length $w$ exceeds aggregation threshold \\
$\mathit{ChargeComposition}(lo,hi)$ & Net charge bounds &
  \PASS{} if $lo \le$ net charge at pH~7.0 $\le hi$ \\
$\mathit{NoLongHydrophobicStretch}(n)$ & Max stretch length &
  \PASS{} if no contiguous hydrophobic stretch $> n$ residues \\
\bottomrule
\end{tabular}
\end{table}

\subsubsection{Immunogenicity Predicates (4)}

Immunogenicity predicates evaluate the risk that the designed protein
will provoke an unwanted immune response, using NetMHCpan (API with
PSSM fallback) for T-cell epitope prediction and B-cell epitope
scoring heuristics.

\begin{table}[t]
\centering
\caption{Immunogenicity type predicates. All are SLOT predicates;
         verdicts depend on NetMHCpan epitope prediction.
         Population-level parameters enable allele-frequency-aware
         assessment.}
\label{tab:imm-predicates}
\begin{tabular}{llp{4.6cm}}
\toprule
\textbf{Predicate} & \textbf{Parameters} & \textbf{Three-valued semantics} \\
\midrule
$\mathit{LowImmunogenicity}(\theta)^\dagger$ & Threshold &
  \PASS{} if immunogenicity score $< \theta$; \UNCERTAIN{} if
  borderline; \FAIL{} if $\ge \theta_{\mathit{hi}}$ \\
$\mathit{NoStrongTCellEpitope}(\theta,\mathit{pop})$ & IC50, population &
  \PASS{} if no peptide binds MHC with IC50 $< \theta$~nM for
  $>$5\% of target population \\
$\mathit{NoDominantBCellEpitope}(\theta)$ & Threshold &
  \PASS{} if no linear B-cell epitope score $\ge \theta$ \\
$\mathit{PopulationCoverageSafe}(p)$ & Coverage threshold &
  \PASS{} if population coverage of strong binders $< p$\% \\
\bottomrule
\end{tabular}
\end{table}

\subsubsection{SLOT vs.\ Core Predicates}

Of the 43~predicates, 17~are \emph{core predicates} (formalized in the
Lean4 model without oracle dependencies) and 19~are \emph{SLOT predicates}
(requiring external oracle calls, formalized via the SLOT framework).
An additional 7~extended diagnostic predicates in the Python
implementation provide heuristic coverage beyond the formal model.
This distinction is critical for the soundness proof:

\begin{itemize}[leftmargin=*,itemsep=0pt,topsep=1pt]
\item \textbf{Core predicates} (17/43) are computationally decidable
      from the sequence alone.
      Their verdicts are \emph{deterministic} and \emph{reproducible}
      by the standalone verifier without any external tool.
      The Lean~4 soundness proof covers these predicates directly.

\item \textbf{SLOT predicates} (19/43) require external oracle input.
      In CONSERVATIVE mode, they always yield \UNCERTAIN{} (the type
      system refuses to guarantee what it cannot verify).
      In VERIFIED mode, they may yield \PASS{} or \FAIL{} based on
      oracle output, but the SLOT-Independence theorem ensures that
      these verdicts are \emph{tracked} separately and do not
      compromise the soundness of core verdicts.
\end{itemize}

\noindent The SLOT framework thus provides a \emph{graduated trust
model}: a CONSERVATIVE-mode certificate's \PASS{} verdicts rest on
only the 17~core predicates (fully verified in Lean~4), while a
VERIFIED-mode certificate additionally trusts the external analysis
engines.
The formal refinement theorem (Refinement.lean, 822~lines,
19~theorems) proves that every \PASS{} verdict in VERIFIED mode is
also valid in CONSERVATIVE mode, establishing that VERIFIED mode
\emph{refines} CONSERVATIVE mode.

\subsection{SLOT: Sealed Local Oracle Theory}\label{sec:slot}

External analysis tools---structure predictors (ESMFold), stability
analyzers (FoldX), solubility scorers (CamSol), immunogenicity
assessors (NetMHCpan)---provide biologically essential information
that cannot be derived from the nucleotide sequence alone.
However, these tools are complex (ESMFold is a 650M-parameter
transformer), potentially unreliable (API outages, version skew), and
certainly not amenable to formal verification in Lean~4.

SLOT provides a principled framework for integrating such tools
without extending the TCB.
The key idea is to treat external tool outputs as \emph{sealed
oracles}: their results are available to the optimizer for guiding
design decisions but are \emph{sealed} (inaccessible) to the type
checker when computing verdicts.

\begin{definition}[SLOT Modes]
\label{def:slot-modes}
BioCompiler operates in one of three SLOT modes:
\begin{description}[leftmargin=!,labelwidth=2.2cm,itemsep=1pt,topsep=1pt]
\item[CONSERVATIVE] No external oracle may be consulted.
      Only core predicates are evaluated; SLOT predicates yield
      \UNCERTAIN{} unconditionally.
      The TCB is minimal: Lean~4 kernel $+$ standalone verifier.

\item[VERIFIED] External oracles may be consulted, but their outputs
      are used only for \FAIL{} verdicts.
      A SLOT predicate may yield \FAIL{} if the oracle definitively
      refutes the property, but may \emph{never} yield \PASS{} based
      on oracle output alone.
      The TCB extends to: ``the oracle correctly identifies
      violations.''

\item[PERMISSIVE] External oracle outputs may influence both \PASS{}
      and \FAIL{} verdicts.
      SLOT predicates may yield \PASS{} based on oracle output.
      The TCB extends to the full correctness of the oracle.
\end{description}
\end{definition}

\begin{theorem}[SLOT Refinement]\label{thm:refinement}
VERIFIED mode refines CONSERVATIVE mode: for every predicate $\pred$,
sequence $\seq$, and context $\Gamma$,
\[
  \judg[\textit{VERIFIED}]{\Gamma}{\seq}{\pred} = \PASS
  \;\;\Longrightarrow\;\;
  \judg[\textit{CONSERVATIVE}]{\Gamma}{\seq}{\pred} = \PASS
\]
\end{theorem}

\begin{proof}[Proof sketch]
In CONSERVATIVE mode, SLOT predicates yield \UNCERTAIN{}, while core
predicates yield the same verdict in both modes (they do not consult
oracles).
The only way $\judg[\textit{VERIFIED}]{\Gamma}{\seq}{\pred} = \PASS$
for a SLOT predicate is if the predicate's condition is satisfied
\emph{without} relying on an oracle \PASS---but VERIFIED mode
prohibits oracle-derived \PASS{} verdicts by definition.
Hence a SLOT predicate can yield \PASS{} in VERIFIED mode only if
its core-computable conditions are satisfied, which are the same
conditions checked in CONSERVATIVE mode.
For core predicates, verdicts are identical across modes.
Since $\sqcap$ preserves the refinement (by absorptivity of \FAIL{}
and the fact that $\UNCERTAIN \sqcap v \neq \PASS$ for any~$v$), the
composite verdict in VERIFIED mode can be \PASS{} only if all
individual verdicts---including the CONSERVATIVE-mode ones---are
\PASS.
\end{proof}

\noindent The refinement theorem is fully formalized and proved in
\texttt{Refinement.lean} (822~lines, 19~theorems, zero
\texttt{sorry}).
It provides the formal foundation for the SLOT trust model: a user who
trusts only the CONSERVATIVE mode guarantees can still benefit from
VERIFIED mode's additional coverage, with the assurance that no
\PASS{} verdict is ``invented'' by the oracle.

\subsection{Typing Rules}\label{sec:typingrules}

The typing judgment has the form
$\judg{\Gamma}{\seq}{\pred} = \verdict$, read: ``under context
$\Gamma$, sequence $\seq$ satisfies predicate $\pred$ with verdict
$\verdict$.''

The structural rules are:

\begin{mathpar}
\inferrule*[Right=T-Pass]
  {\mathit{propertyHolds}(\pred, \seq, \Gamma) \\
   \mathit{certainly}(\pred, \seq, \Gamma)}
  {\judg{\Gamma}{\seq}{\pred} = \PASS}
\and
\inferrule*[Right=T-Fail]
  {\mathit{propertyRefuted}(\pred, \seq, \Gamma)}
  {\judg{\Gamma}{\seq}{\pred} = \FAIL}
\and
\inferrule*[Right=T-Uncertain]
  {\neg\, \mathit{certainly}(\pred, \seq, \Gamma) \\
   \neg\, \mathit{propertyRefuted}(\pred, \seq, \Gamma)}
  {\judg{\Gamma}{\seq}{\pred} = \UNCERTAIN}
\end{mathpar}

\noindent The side conditions are:
\begin{itemize}[leftmargin=*,itemsep=0pt,topsep=1pt]
\item $\mathit{propertyHolds}(\pred, \seq, \Gamma)$: the property
      named by $\pred$ is true of $\seq$ in context $\Gamma$.
\item $\mathit{certainly}(\pred, \seq, \Gamma)$: the evidence for
      the property is sufficient to establish a guarantee (e.g., the
      score is below the uncertain threshold $\theta_u$, not merely
      below the failure threshold $\theta_c$).
\item $\mathit{propertyRefuted}(\pred, \seq, \Gamma)$: the property
      is provably violated (e.g., a restriction site is present, or
      a splice site score exceeds $\theta_c$).
\end{itemize}

\noindent Composition via Kleene conjunction:

\begin{mathpar}
\inferrule*[Right=T-And]
  {\judg{\Gamma}{\seq}{\pred_1} = v_1 \\
   \judg{\Gamma}{\seq}{\pred_2} = v_2 \\
   v = v_1 \sqcap v_2}
  {\judg{\Gamma}{\seq}{\pred_1 \sqcap \pred_2} = v}
\end{mathpar}

\noindent Per-predicate rules specialize
\textsc{T-Pass}/\textsc{T-Fail}/\textsc{T-Uncertain} to each
biological property.
We show four representative rules, one from each major category:

\paragraph{DNA-level: NoRestrictionSite.}
This predicate checks that no restriction enzyme recognition sequence
from set $S$ occurs in $\seq$, including at codon boundaries and
fragment junctions.
It is a core predicate (no SLOT dependency) and uses exact string
matching.

\begin{mathpar}
\inferrule*[Right=T-Restr-Pass]
  {\forall\, s \in S.\; \neg\,\mathit{containsPattern}(\seq, s)}
  {\judg{\Gamma}{\seq}{\mathit{NoRestrictionSite}(S)} = \PASS}
\and
\inferrule*[Right=T-Restr-Fail]
  {\exists\, s \in S.\; \mathit{containsPattern}(\seq, s)}
  {\judg{\Gamma}{\seq}{\mathit{NoRestrictionSite}(S)} = \FAIL}
\end{mathpar}

\noindent Note that \textsc{NoRestrictionSite} is purely binary
(\PASS{} or \FAIL); there is no \UNCERTAIN{} case because pattern
matching is exact.
This is characteristic of predicates based on exact computation:
uncertainty arises only when the predicate involves a scoring function
with continuous output.

\paragraph{DNA-level with dual threshold: NoCrypticSplice.}
This predicate uses MaxEntScan scores with dual thresholds
$\theta_u$ and $\theta_c$, producing all three verdicts.

\begin{mathpar}
\inferrule*[Right=T-Cryptic-Pass]
  {\forall\, p \in \mathit{spliceSites}(\seq).\;
   \mathit{score}(p) < \theta_u}
  {\judg{\Gamma}{\seq}{\mathit{NoCrypticSplice}(\theta_c,\theta_u)} = \PASS}
\and
\inferrule*[Right=T-Cryptic-Uncertain]
  {\forall\, p.\; \mathit{score}(p) < \theta_c \\
   \exists\, p \in \mathit{spliceSites}(\seq).\;
   \theta_u \le \mathit{score}(p) < \theta_c}
  {\judg{\Gamma}{\seq}{\mathit{NoCrypticSplice}(\theta_c,\theta_u)} = \UNCERTAIN}
\and
\inferrule*[Right=T-Cryptic-Fail]
  {\exists\, p \in \mathit{spliceSites}(\seq).\;
   \mathit{score}(p) \ge \theta_c}
  {\judg{\Gamma}{\seq}{\mathit{NoCrypticSplice}(\theta_c,\theta_u)} = \FAIL}
\end{mathpar}

\noindent The \textsc{T-Cryptic-Uncertain} rule requires both
\emph{absence} of any definitively failing site (all scores $<
\theta_c$) and \emph{presence} of at least one borderline site.
This ensures that \UNCERTAIN{} does not mask a \FAIL: if any site
exceeds $\theta_c$, the verdict is \FAIL{} regardless of borderline
sites.

\paragraph{Structure SLOT predicate: StructureConfidence.}
This SLOT predicate relies on ESMFold for structure prediction.
In CONSERVATIVE mode, it always yields \UNCERTAIN{}.
In VERIFIED mode, the oracle may establish \FAIL{} but not \PASS.

\begin{mathpar}
\inferrule*[Right=T-Struct-Cons]
  {\mathit{mode} = \textit{CONSERVATIVE}}
  {\judg{\Gamma}{\seq}{\mathit{StructureConfidence}(\theta)} = \UNCERTAIN}
\and
\inferrule*[Right=T-Struct-Fail]
  {\mathit{mode} \in \{\textit{VERIFIED}, \textit{PERMISSIVE}\} \\
   \mathit{oracle\_pLDDT}(\seq) < \theta_{\mathit{lo}}}
  {\judg{\Gamma}{\seq}{\mathit{StructureConfidence}(\theta)} = \FAIL}
\and
\inferrule*[Right=T-Struct-Pass]
  {\mathit{mode} = \textit{PERMISSIVE} \\
   \mathit{oracle\_pLDDT}(\seq) \ge \theta}
  {\judg{\Gamma}{\seq}{\mathit{StructureConfidence}(\theta)} = \PASS}
\end{mathpar}

\noindent The \textsc{T-Struct-Cons} rule is unconditional in
CONSERVATIVE mode: the type system refuses to evaluate structural
properties without oracle input, yielding \UNCERTAIN{} by default.
The \textsc{T-Struct-Fail} rule allows the oracle to establish a
\FAIL{} verdict in VERIFIED mode---this is sound because a \FAIL{}
verdict does not compromise the soundness guarantee (soundness
constrains \PASS{}, not \FAIL).
The \textsc{T-Struct-Pass} rule is available only in PERMISSIVE mode,
where the oracle is fully trusted.

\paragraph{Immunogenicity SLOT predicate: NoStrongTCellEpitope.}
This SLOT predicate relies on NetMHCpan for peptide--MHC binding
prediction.
It is parameterized by an IC50 threshold and a target population
(specified as a set of HLA alleles with frequencies).

\begin{mathpar}
\inferrule*[Right=T-TCR-Pass]
  {\mathit{mode} = \textit{PERMISSIVE} \\
   \forall\, \mathit{pep} \in \mathit{peptides}(\seq).\;
   \forall\, h \in \mathit{pop}.\;
   \mathit{IC50}(\mathit{pep}, h) \ge \theta}
  {\judg{\Gamma}{\seq}{\mathit{NoStrongTCellEpitope}(\theta, \mathit{pop})} = \PASS}
\and
\inferrule*[Right=T-TCR-Fail]
  {\mathit{mode} \in \{\textit{VERIFIED}, \textit{PERMISSIVE}\} \\
   \exists\, \mathit{pep}, h.\;
   \mathit{IC50}(\mathit{pep}, h) < \theta \;\wedge\;
   \mathit{freq}(h) > 0.05}
  {\judg{\Gamma}{\seq}{\mathit{NoStrongTCellEpitope}(\theta, \mathit{pop})} = \FAIL}
\and
\inferrule*[Right=T-TCR-Cons]
  {\mathit{mode} = \textit{CONSERVATIVE}}
  {\judg{\Gamma}{\seq}{\mathit{NoStrongTCellEpitope}(\theta, \mathit{pop})} = \UNCERTAIN}
\end{mathpar}

\noindent The population-frequency threshold ($>5\%$) ensures that
rare-allele binders do not trigger a \FAIL{} verdict, reflecting
clinical practice where population-level immunogenicity assessment
focuses on common HLA alleles.
The fallback to PSSM-based prediction when the NetMHCpan API is
unavailable provides graceful degradation: the PSSM predictor is less
accurate but produces conservative (toward \FAIL{}) estimates,
maintaining soundness.

\subsection{Soundness and Compositionality}%
\label{sec:compositionality}

The central guarantee of BioCompiler's type system is that \PASS{}
verdicts are never wrong.
This is formalized as a type soundness theorem and a compositional
soundness corollary, both proved in Lean~4.

\begin{theorem}[Type Soundness]\label{thm:sound}
For all predicates $\pred$, sequences $\seq$, and contexts $\Gamma$:
\[
  \judg{\Gamma}{\seq}{\pred} = \PASS
  \;\longrightarrow\;
  \mathit{propertyHolds}(\pred, \seq, \Gamma)
\]
\end{theorem}

\begin{proof}[Proof sketch]
By case analysis on the predicate~$\pred$.
For core predicates (e.g., $\mathit{NoRestrictionSite}$,
$\mathit{GCInRange}$, $\mathit{InFrame}$), the \PASS{} branch of
the evaluation function directly checks the property (e.g., the
scanner finds no restriction site, the GC computation falls in range),
so soundness follows by computation.
For dual-threshold predicates (e.g., $\mathit{NoCrypticSplice}$),
the \PASS{} branch requires all scores $< \theta_u$, which is a
stronger condition than ``no score $\ge \theta_c$''; the gap between
$\theta_u$ and $\theta_c$ provides a safety margin.
For SLOT predicates in CONSERVATIVE mode, \PASS{} is unreachable
(the verdict is always \UNCERTAIN{}), so the implication is vacuously
true.
For SLOT predicates in VERIFIED mode, \PASS{} is reachable only via
core-computable conditions (by the VERIFIED-mode restriction), so
soundness reduces to the core-predicate case.
The Lean~4 proof (module \texttt{TypeSystem.lean}) performs this case
analysis for all 36~predicates in the Lean4 model.
\end{proof}

\begin{remark}[Vacuous soundness of SLOT predicates]\label{rem:vacuous-slot}
In the current Lean4 formalization, all 19 SLOT-dependent predicates
evaluate to \UNCERTAIN{} regardless of input, making the per-predicate
soundness theorem for SLOT predicates vacuously true: since \PASS{} is
never produced, the implication $\PASS \to \mathit{propertyHolds}$ holds
trivially. This is by design---the \CONSERVATIVE{} mode guarantees that
no SLOT predicate can contribute to a \PASS{} verdict, ensuring that the
soundness guarantee depends only on the verified core predicates. The
\texttt{SLOTVerification} module introduces verification conditions that,
when discharged, allow SLOT predicates to return non-\UNCERTAIN{} verdicts
in \VERIFIED{} mode (implemented in Python but not yet formally verified).
The real guarantee for SLOT predicates thus comes from the Python
implementation in \VERIFIED{} mode, which is not formally verified; the
Lean4 model is sound by construction, and the \texttt{SLOTVerification}
module provides the framework for progressive assurance as oracles are
validated.
\end{remark}

\begin{corollary}[Compositional Soundness]\label{cor:comp}
If\/ $\judg{\Gamma}{\seq}{\pred_1 \sqcap \pred_2} = \PASS$, then
$\mathit{propertyHolds}(\pred_1, \seq, \Gamma)$ and
$\mathit{propertyHolds}(\pred_2, \seq, \Gamma)$.
\end{corollary}

\begin{proof}
Since $\judg{\Gamma}{\seq}{\pred_1 \sqcap \pred_2} = \PASS$ and
$\sqcap$ is soundness-preserving ($v_1 \sqcap v_2 = \PASS \Rightarrow
v_1 = v_2 = \PASS$), both individual verdicts must be \PASS.
By Theorem~\ref{thm:sound}, each predicate's property holds.
\end{proof}

\begin{corollary}[SLOT-Independence]\label{cor:slot}
In CONSERVATIVE mode, FFI output does not affect verdicts:
$\judg[\textit{CONSERVATIVE}]{\Gamma}{\seq}{\pred}$ depends only on
non-FFI state.
\end{corollary}

\begin{proof}
In CONSERVATIVE mode, all SLOT predicates yield \UNCERTAIN{}
unconditionally (by \textsc{T-Struct-Cons} and analogous rules for
other SLOT predicates), so their verdicts are constant and do not
depend on oracle output.
Core predicates do not consult oracles by definition.
\end{proof}

\noindent The compositional soundness corollary is the formal
guarantee that addresses the junction counterexample of
\S\ref{sec:motivating}: when predicates are combined via $\sqcap$,
a composite \PASS{} verdict guarantees that \emph{every individual
property holds on the entire sequence}, including junction regions
that are invisible to per-fragment checking.
This is precisely what constraint lists cannot provide: there is no
analogous theorem for the conjunction of independent constraint checks,
because the checks are performed on \emph{fragments}, not on the
\emph{composed artifact}.

The soundness proof is fully machine-verified in Lean~4 across
17~modules (12~proof $+$ 5~library), totaling 8,495~lines of proof
code with zero \texttt{sorry} in the core type system and zero standalone
\texttt{axiom} declarations; 3 class-field axioms in SpliceSiteScanner
(scanner\_completeness, scanner\_soundness, borderline\_completeness)
remain as TCB parameters, and 15 proof obligations in the SLOT verification
theorem use \texttt{sorry} as placeholders.
This represents a significant scaling effort from the original 8-predicate,
2,400-line proof: the 43-predicate system required restructuring the
proof architecture into a layered dependency graph
(Figure~\ref{fig:proofdep}), with per-domain soundness modules that
feed into the central compositionality theorem.

\begin{figure}[t]
\centering
\begin{tikzpicture}[
  node distance=0.45cm and 0.9cm,
  box/.style={draw,rounded corners,minimum height=0.65cm,minimum
              width=2.4cm,font=\scriptsize,align=center,fill=white},
  dom/.style={draw,rounded corners,minimum height=0.65cm,minimum
              width=2.4cm,font=\scriptsize,align=center,fill=blue!8},
  arr/.style={-{Stealth[length=4pt]},thick}
]
\node[box] (tv) {ThreeValued\\(K3 algebra)};
\node[box,below=of tv] (seq) {Sequence\\(scanner)};
\node[dom,right=of tv] (dna) {DNA-level\\(12 predicates)};
\node[dom,right=of seq] (slot) {SLOT\\Framework};
\node[dom,below=0.7cm of dna] (struct) {Structure\\(4 predicates)};
\node[dom,right=of struct] (stab) {Stability\\(4 predicates)};
\node[dom,below=0.7cm of slot] (sol) {Solubility\\(4 predicates)};
\node[dom,right=of sol] (imm) {Immunogenicity\\(4 predicates)};
\node[box,below=1.0cm of stab] (comp) {Compositional\\Soundness};
\node[box,below=1.0cm of imm] (ref) {Refinement\\(SLOT)};

\draw[arr] (tv) -- (dna);
\draw[arr] (seq) -- (dna);
\draw[arr] (tv) -- (slot);
\draw[arr] (dna) -- (struct);
\draw[arr] (dna) -- (stab);
\draw[arr] (slot) -- (struct);
\draw[arr] (slot) -- (stab);
\draw[arr] (slot) -- (sol);
\draw[arr] (slot) -- (imm);
\draw[arr] (struct) -- (comp);
\draw[arr] (stab) -- (comp);
\draw[arr] (sol) -- (comp);
\draw[arr] (imm) -- (comp);
\draw[arr] (slot) -- (ref);
\draw[arr] (comp) -- (ref);
\end{tikzpicture}
\caption{Proof dependency graph for the 43-predicate type system.
         Blue nodes are per-domain soundness modules; white nodes are
         core libraries. The graph shows how the original flat
         architecture was restructured into a layered design to
         scale from 8 to 43 predicates.}
\label{fig:proofdep}
\end{figure}

The proof architecture (Figure~\ref{fig:proofdep}) reflects the
domain structure of the predicates: per-domain soundness modules
(DNA-level, Structure, Stability, Solubility, Immunogenicity) each
prove that their predicates are sound individually, then feed into
the central Compositional Soundness module.
The SLOT Refinement module is orthogonal: it proves that VERIFIED
mode refines CONSERVATIVE mode across all domains, using the
per-domain soundness results as lemmas.

This modular architecture was essential for scaling the proof from 8
to 43 predicates without exponential blowup.
Each domain is proved independently, and the compositionality theorem
stitches them together using the algebraic properties of Kleene
conjunction.
Adding a new predicate to an existing domain requires updating only
that domain's soundness module and the compositionality theorem's
case split---a local change, not a global refactoring.
